\documentclass[12pt,a4paper,oneside]{article}
\usepackage[ngerman]{babel}
\usepackage[utf8]{inputenc}
\usepackage[babel,style=german]{csquotes}
\usepackage{amsmath}
\usepackage{amssymb}
\usepackage{amsthm}
\usepackage{booktabs}
\usepackage{changepage}
\usepackage{enumitem}
\usepackage{xcolor}
\usepackage{empheq}
\usepackage{fancyhdr}
\usepackage{graphicx}
\usepackage{hyperref}
% \usepackage[draft]{hyperref}
\usepackage{icomma}
\usepackage{soul}
\setul{1.2pt}{0.5pt}
\usepackage{tikz}
\usetikzlibrary{shapes.multipart}
\usepackage{tkz-euclide}
\usepackage{titlesec}
\usepackage{tabularray}
\UseTblrLibrary{booktabs}
\usepackage{xurl}
\usepackage{geometry}
\geometry{
    left=2.28cm,
    right=2.28cm,
    top=3cm,
    bottom=3cm
}

\setlength{\parindent}{0pt}

\definecolor{myorange}{HTML}{B35900}
\definecolor{mygreen}{HTML}{009900}
\definecolor{mydarkgreen}{HTML}{005500}
\definecolor{mybrightyellow}{RGB}{255, 243, 205}

\newcommand{\afz}[1]{\enquote{#1}}
\newcommand{\emptyline}{\hspace*{\fill}}
\newcommand{\N}{\mathbb{N}}
\newcommand{\R}{\mathbb{R}}
\newcommand{\If}[1]{\text{falls}\; #1}
\newcommand{\ch}[3]{%
    \begin{tabular}[b]{@{\hspace{2pt}}c@{\hspace{2pt}}}
        \textcolor{#1}{#2} \\[2pt]
        \textcolor{blue}{\texttt{#3}}
    \end{tabular}%
}

\begin{document}

    \tableofcontents

    \newpage

    \section{Einführung}

        Es gibt unglaublich viele Sprachen.
        Aber wenn wir alle Sprachen abstrahieren, können wir erkennen, dass sie aus Zeichen, einem Alphabet und Wörtern gebildet werden.

        \emptyline

        Das Zeichen ist die kleinste Einheit einer Sprache.
        Wir können die Zeichen als Symbole verstehen.
        Alle erlaubten Zeichen bilden ein Alphabet, als eine endliche Menge von Zeichen.
        Und daraus können wir Wort definieren: Ein Wort ist eine endliche Zeichenfolge.
        Ähnlicheweise können wir auch Sprache definieren: Eine Sprache ist eine Menge von Wörtern, die alle erlaubten Wörter enthält.
        Es ist leicht zu sehen, dass Sprachen meistens unendlich groß sind, da es normalerweise unendlich viele erlaubte Wörter gibt,
        die wir uns nicht alle einzeln vorstellen können.

    \section{Alphabet}

        Ein Alphabet ist mathematisch eine endliche, nicht-leere Menge.
        Es wird oftmals mit $A$ bezeichnet, aber später wird auch $\Sigma$ oft verwendet.
        Die Elemente dieser Menge nennen wir entweder Zeichen oder Symbole.

        \emptyline

        Wir stellen einige Beispiele in der Theoretischen Informatik vor:

        \vspace{-1em}
        \begin{align*}
            &A = \{0, 1\} \\
            &A = [n] = \{i \in \N_+ \;|\; 1 \leq i \leq n\} \; (n \in \N_+) \\
            &A = \{!, *, \%\}
        \end{align*}

        Die Beispiele sind nicht sehr anschaulich, deswegen nennen wir auch paar Beispiele in der Praktischen Informatik:

        \begin{itemize}
            \item ASCII
            \item Unicode
            \item ARMSCII
            \item DIN 91379
            \item ...
        \end{itemize}

    \newpage

    \section{Wörter}

        \subsection{Grundlagen}

            Wir nehmen an, dass ein Alphabet $A$ fest ist.
            Es wird einmal fixiert, und nicht zwischendrin verändert.

            \emptyline

            Dann definieren wir ein Wort als eine endliche Folge von Zeichen aus $A$.
            Zu beachten ist, dass die Symbole mehrfach vorkommen dürfen.

            \emptyline

            Ein Wort (über dem Alphabet $A$) kann nummeriert werden.
            Die Zeichen-Reihenfolge nennen wir Nummerierung.
            Das heißt, dass ein Wort tatsächlich eine Abbildung $w: [n] \to A$ ist.

            \begin{center}
                \begin{tblr}{
                    colspec = {*{9}{c}},
                    row{1} = {fg=myorange},
                    row{2} = {fg=blue, font=\ttfamily}
                }
                    \toprule
                    1 & 2 & 3 & 4 & 5 & 6 & 7 & 8 & 9 \\
                    M & i & l & c & h & r & e & i & s \\
                    \bottomrule
                \end{tblr}
            \end{center}

            $n$ heißt die Länge eines Wortes, geschrieben $|w|$.
            normalerweise gilt immer $n \geq 1$.
            Die einzige Ausnahme ist das \textit{leere} Wort $\varepsilon$ (kommt später dazu).

            \emptyline

            Wir stellen ein Beispiel \afz{hallo} vor:

            \vspace{-1em}
            \begin{gather*}
                w : [5] \to \{\text{a}, \text{h}, \text{l}, \text{o}\}, \\
                w(1) = \text{h},\ w(2) = \text{a},\ w(3) = w(4) = \text{l},\ w(5) = \text{o}
            \end{gather*}

            bzw.:

            \vspace{-1em}
            \begin{align*}
                1 \mapsto \text{h},\ 2 \mapsto \text{a},\ 3 \mapsto \text{l},\ 4 \mapsto \text{l},\ 5 \mapsto \text{o}
            \end{align*}

            Um es einfacher zu lesen, schreiben wir normalerweise $w_i$ statt $w(i)$.

            \emptyline

            Wir definieren $A^*$ als die Menge aller möglicher Wörter über Alphabet $A$.
            Das heißt, alle Wörter, die nur Zeichen aus $A$ enthalten, sind Elemente von $A^*$.

            \emptyline

            Zum Beispiel sei $A = \{\text{a}, \text{b}\}$, dann enthält $A^*$ die Wörter:

            \begin{itemize}
                \item a und b
                \item aa, ab, ba und bb
                \item aaa, aab, aba, abb, baa, bab, bba und bbb,
                \item ...
                \item und außerdem $\varepsilon$, das \textit{leere} Wort.
            \end{itemize}

            Wir können leicht erkennen, dass es unendlich viele Wörter in $A^*$ gibt, die alle endliche Länge haben.
            Dadurch definieren wir $A^*$ formal:

            \vspace{-1em}
            \begin{align*}
                A^* = \{w: [n] \to A \;|\; n \in \N_0\}
            \end{align*}

        \subsection{Das leere Wort}

            Das leere Wort besteht aus 0 Symbolen.
            Notiert als $\varepsilon$.

            \emptyline

            Das leere Wort ist eindeutig, weil:

            \vspace{-1em}
            \begin{align*}
                \varepsilon: [0] \to A \Longleftrightarrow \varepsilon: \emptyset \to A
            \end{align*}

            Diese Abbildung hat einen leeren Definitionsbereich, deswegen gibt es nur eine solche Abbildung.
            Daraus folgt, dass das leere Wort eindeutig ist.

            \emptyline

            \textbf{ACHTUNG} das leere Wort ist \afz{etwas}.
            Die Kardinalität der Menge $\{\varepsilon\}$ ist $|\{\varepsilon\}| = 1$.
            Das ist nicht die leere Menge (Zum Vergleich $|\emptyset| = 0$).

        \subsection{\texorpdfstring{$A^n$ - die Wörter der Länge $n$}{Aⁿ - die Wörter der Länge n}}

            Wir definieren $A^n$ als die Menge aller Wörter der Länge $n \in \N_0$ über dem Alphabet $A$.

            \emptyline

            Also ist sozusagen:

            \vspace{-1em}
            \begin{align*}
                A^* = A^0 \cup A^1 \cup A^2 \cup A^3 \cup \cdots
            \end{align*}

            Wir können ohne Pünktchen schreiben, damit die Definition exakt ist:

            \vspace{-1em}
            \begin{align*}
                A^* = \bigcup_{i \in \N_0} A^i
            \end{align*}

        \subsection{Konkatenation von Wörtern}

            Wir können zwei Wörter miteinander verketten.
            Diese Operation wird als Konkatenation bezeichnet und durch das Symbol \afz{$\cdot$} dargestellt.

            \emptyline

            Eine formale Definition ist:
            Es seien $w_1: [m] \to A_1$ und $w_2: [n] \to A_2$ zwei Wörter.
            Dann gilt:

            \vspace{-1em}
            \begin{align*}
                w_1 \cdot w_2: [m + n] &\to A_1 \cup A_2 \\
                                     i &\mapsto
                \begin{cases}
                    w_1(i)     &\If{1 \leq i \leq m} \\
                    w_2(i - m) &\If{m < i \leq m + n}
                \end{cases}
            \end{align*}

            Übrigens gilt:

            \vspace{-1em}
            \begin{align*}
                |w_1 \cdot w_2| = |w_1| + |w_2|
            \end{align*}

            Ein Beispiel stellen wir vor:

            \vspace{-1em}
            \begin{align*}
                &\ch{myorange}{1}{A}\ \ch{myorange}{2}{P}\ \ch{myorange}{3}{F}\ \ch{myorange}{4}{E}\ \ch{myorange}{5}{L}
                \ \cdot \ 
                \ch{mygreen}{1}{M}\ \ch{mygreen}{2}{U}\ \ch{mygreen}{3}{S} \\[10pt]
                =\ &\ch{myorange}{1}{A}\ \ch{myorange}{2}{P}\ \ch{myorange}{3}{F}\ \ch{myorange}{4}{E}\ \ch{myorange}{5}{L}\ \ch{mydarkgreen}{6}{M}\ \ch{mydarkgreen}{7}{U}\ \ch{mydarkgreen}{8}{S}
            \end{align*}

            Den Punkt können wir manchmal weglassen.

            \newpage

            Für jedes Alphabet $A$ gilt:

            \vspace{-1em}
            \begin{align*}
                \forall w \in A^*: w \cdot \varepsilon = w \wedge \varepsilon \cdot w = w
            \end{align*}

            Das ist anschaulich klar, aber wir können es natürlich auch formal beweisen.

            \emptyline

            \begin{adjustwidth}{2em}{0pt}
                Sei $w' := w \cdot \varepsilon$ mit $|w| = m$.
                Dann ist $w'$ eine Abbildung $w': [m + 0] \to A$.

                \emptyline

                Dann für $i \in [m]$ gilt:

                \vspace{-1em}
                \begin{align*}
                    w'(i) = (w \cdot \varepsilon)(i) &=
                    \begin{cases}
                                      w(i) &\If{1 \leq i \leq m} \\
                        \varepsilon(i - m) &\If{m < i \leq m + 0}
                    \end{cases} \\
                                                     &= w(i)
                \end{align*}

                Und es ist leicht zu sehen, dass $w$ und $w'$ gleichen Definitionsbereich und Zielbereich haben.
                Gleichzeitig haben sie überall die gleichen Funktionswerte.
                Also $w' = w$, folgt $w \cdot \varepsilon = w$.
                Den umgekehrten Fall $\varepsilon \cdot w = w$ können wir auf dieselbe Weise beweisen.
            \end{adjustwidth}

            \emptyline

            Die Konkatenation von Wörtern ist \textbf{nicht kommutativ}.
            Das heißt, dass die Reihenfolge wichtig ist.
            Aber Konkatenation ist assoziativ, was heißt (hier sind $w_1$, $w_2$ und $w_3$ Wörtern):

            \vspace{-1em}
            \begin{align*}
                \forall w_1, w_2, w_3 \in A^*: (w_1 \cdot w_2) \cdot w_3 = w_1 \cdot (w_2 \cdot w_3)
            \end{align*}

            Deswegen ist die Klammerung egal.

            \emptyline

            Analog zur Potenzschreibweise bei Zahlen definieren wir für Wörter ähnlicheweise wie

            \vspace{-1em}
            \begin{align*}
                w^n = \underbrace{w \cdot w \cdot \cdots \cdot w}_{n \;\text{mal}}
            \end{align*}

            Aber da wir Notation mit Pünktchen vermeiden müssen, schreiben wir wie:

            \vspace{-1em}
            \begin{align*}
                \forall n \in \N_0\ \forall w \in A^*: w^n =
                \begin{cases}
                    \varepsilon     &\If{n = 0} \\
                    w \cdot w^{n-1} &\If{n \geq 1}
                \end{cases}
            \end{align*}
        
        \subsection{Teilwörter}

            Ein Teilwort (= Substring) ist eine zusammenhängende Teilfolge.

            \emptyline

            Seien $v, w \in A^*$ Wörter über einem Alphabet $A$, dann ist:

            \begin{itemize}
                \item $v$ ein Teilwort von $w$ wenn:
                    \begin{align*}
                        \exists x_1, x_2 \in A^*: w = x_1 \cdot v \cdot x_2
                    \end{align*}
                \item $v$ ein Präfix von $w$ wenn:
                    \begin{align*}
                        \exists x \in A^*: w = v \cdot x
                    \end{align*}
                \item $v$ ein Suffix von $w$ wenn:
                    \begin{align*}
                        \exists x \in A^*: w = x \cdot v
                    \end{align*}
            \end{itemize}

    \section{Sprachen}

        \subsection{Grundlagen}

            Es gibt natürliche Sprache und künstliche Sprachen.
            Man fragt sich, wie man spezifiziert, was syntaktisch korrekt ist.
            Man fragt sich auch, wie man überprüft, ob etwas syntaktisch korrekt ist.
            Wie definiert man die Bedeutung syntaktisch korrekter Gebilde?

            \emptyline

            Sei $A$ ein Alphabet.
            Dann ist eine \textit{formale Sprache} $L$ eine Menge von Wörtern über $A$, also $L \subseteq A^*$.
            Hier spielt die Bedeutung keine Rolle.

            \emptyline

            Formale Sprachen sind als Modelle / Formalisierungen sehr vielseitig.
            
            \begin{itemize}
                \item für Syntax: \afz{Wörter der korrekten Form}.
                \item für Semantik: \afz{Wörter der korrekten Bedeutung}.
            \end{itemize}

        \subsection{Produkte formaler Sprachen}

            Sei zwei formale Sprachen $L_1$ und $L_2$.
            Dann ist das Produkt der zwei Sprachen:

            \vspace{-1em}
            \begin{align*}
                L_1 \cdot L_2 = \{w_1 \cdot w_2 \;|\; w_1 \in L_1 \wedge w_2 \in L_2\}
            \end{align*}

            Wegen der Assoziativität der Konkatenation von Wörtern folgt:

            \vspace{-1em}
            \begin{align*}
                L_1 \cdot L_2 \cdot L_3 = \{w_1 \cdot w_2 \cdot w_3 \;|\; w_1 \in L_1 \wedge w_2 \in L_2 \wedge w_3 \in L_3\}
            \end{align*}

            Wie bei der Konkatenation von Wörtern lassen wir auch gerne den Punkt wieder weg.

            \emptyline

            Es ist einfach zu beweisen, dass die leere Sprache ein neutrales Element für Produkt formaler Sprachen ist,
            was heißt:

            \vspace{-1em}
            \begin{align*}
                L \cdot \{\varepsilon\} = \{\varepsilon\} \cdot L = L
            \end{align*}

            Der Beweis ist folgende:

            \begin{adjustwidth}{2em}{0pt}
                \vspace{-1em}
                \begin{align*}
                    L \cdot \{\varepsilon\} &= \{w_1w_2 \;|\; w_1 \in L \wedge w_2 \in \{\varepsilon\}\} \\
                                            &= \{w_1w_2 \;|\; w_1 \in L \wedge w_2 = \varepsilon\} \\
                                            &= \{w_1\varepsilon \;|\; w_1 \in L\} \\
                                            &= \{w_1 \;|\; w_1 \in L\} \\
                                            &= L
                \end{align*}

                Analog kann man $L = \{\varepsilon\} \cdot L$ zeigen.
            \end{adjustwidth}

            \emptyline

            Folglich wird die Potenz einer Sprache $L$ definiert: $L^n$ enthält alle Konkatenationen von $n$ Wörtern in $L$:

            \vspace{-1em}
            \begin{align*}
                                    L^0 &= \{\varepsilon\} \\
                \forall n \in \N_+: L^n &= L \cdot L^{n-1}
            \end{align*}

            Genau wie bei Wörtern ist die Konkatenation von Sprachen auch eine assoziative Operation.

            \newpage

            Jedes Wort $w \in L^n$ ist ein Wort $w = w_1 \cdot \cdots \cdot w_n$ mit $w_i \in L$ für jedes $i \in [n]$.

            \emptyline

            \textbf{ABER} verschiedene Potenzen der Sprache $L$ sind nicht unbedingt disjunkt.
            Zum Beispiel seien:

            \vspace{-1em}
            \begin{align*}
                A &= \{\text{a}, \text{b}\} \\
                L &= \{\text{a}, \text{ab}, \text{b}, \text{ba}\}
            \end{align*}

            Dann gibt es (\afz{$|$} als Trennzeichen):

            \vspace{-1em}
            \begin{align*}
                w_1 \in L^2 &= \text{ab|ab} \\
                w_2 \in L^3 &= \text{ab|a|b}
            \end{align*}

            Daraus folgt, dass die Vereinigung von $L^i$ und $L^j$ mit $i, j \in \N_+, i \neq j$ nicht unbedingt leer ist.

        \subsection{Konkatenationsabschluss formaler Sprachen}

            Bei Alphabeten haben wir definiert:

            \vspace{-1em}
            \begin{align*}
                A^* = \bigcup_{i \in \N_0} A^i
            \end{align*}

            Ähnlicheweise definieren wir den Konkatenationsabschluss $L^*$ von $L$ (kleenescher Abschluss):

            \vspace{-1em}
            \begin{align*}
                L^* = \bigcup_{i \in \N_0} L^i
            \end{align*}

            Gleichzeitig gibt es den $\varepsilon$-freier Konkatenationsabschluss $L^+$ von $L$ (positiver Abschluss):

            \vspace{-1em}
            \begin{align*}
                L^+ = \bigcup_{i \in \N_+} L^i
            \end{align*}

            Daraus folgt:

            \vspace{-1em}
            \begin{align*}
                L^* = L^0 \cup L^+
            \end{align*}

            Da jedes Wort aus $A^*$ die Konkatenation von Blöcken, und $L^*$ enthält alle mögliche Konkatenationen von Wörtern,
            gilt $L^* = A^*$.

            \emptyline

            Es ist auch sehr wichtig zu wissen, dass der positive Abschluss nicht \afz{$\varepsilon$-frei} sein muss.
            Da $L = L^i \subseteq L^+$, wenn $\varepsilon \in L$ gilt, dann gilt auch $\varepsilon \in L^+$.
            
            \emptyline

            Am Ende ergänzen wir:

            \vspace{-1em}
            \begin{align*}
                \emptyset^* = \{\varepsilon\}
            \end{align*}

            Später wird es nützlich.

        \newpage

        \subsection{Reguläre Ausdrücke}

            Seien $L_1$, $L_2$ zwei Sprachen.
            Durch verschiedene Operationen wie Produkt, kleenescher Abschluss und Vereinigung entstehen neue Sprachen.
            Wir fragen uns, welche Sprachen so erzeugt werden können, welche davon nicht erzeugt werden können?

            \emptyline

            Sei Alphabet $Z = \{\cup, (, ), *, \emptyset\}$ von \afz{Hilfssymbolen}.
            Alphabet $A$ darf kein Zeichen aus $Z$ enthalten.
            Dann heißt ein regulärer Ausdruck über $A$ ein Wort über dem Alphabet $A \cup Z$,
            das gewissen Vorschriften genügt.

            \emptyline

            Die Menge der regulären Ausdrücke (RA) ist wie folgt festgelegt:

            \begin{itemize}
                \item $\emptyset$ ist ein RA
                \item für jedes Zeichen $a \in A$ ist $a$ ein RA
                \item wenn $R_1$ und $R_2$ RA sind, dann auch $(R_1 \cup R_2)$ und $(R_1 \cdot R_2)$
                \item wenn $R$ ein RA ist, dann auch $(R^*)$
                \item Nichts anderes sind reguläre Ausdrücke
            \end{itemize}

            Das ergibt eine Sprache $L_\text{RA} \subseteq (A \cup Z)^*$.

            \emptyline

            Es gibt eine Regeln, die Klammern zu einsparen.

            \begin{itemize}
                \item äußere Klammern kann man weglassen
                \item \afz{Stern- vor Punktrechnung}
                \item \afz{Punkt- vor Strichrechnung}
                \item bei mehreren gleichen binären Operatoren ohne Klammern gilt das als links geklammert
            \end{itemize}

            Es gibt ein paar Beispiele:

            \vspace{-1em}
            \begin{align*}
                                 R_1 \cup R_2 R_3^* &= R_1 \cup R_2 (R_3^*) \\
                                                    &= R_1 \cup (R_2 (R_3^*)) \\
                                                    &= (R_1 \cup (R_2 (R_3^*))) \\
                              R_1 \cup R_2 \cup R_3 &= ((R_1 \cup R_2) \cup R_3) \\ 
                (R_1 R_2 R_3)^{**} \cup \emptyset^* &= (((((R_1 R_2)R_3)^*)^*) \cup (\emptyset^*))
            \end{align*}

            Wir können eine formale Sprache $L$ durch einen regulären Ausdruck $R$ beschreiben, notiert als $L(R)$.
            Die Definition der Sprache $L(R)$ folgt der des regulären Ausdrucks $R$:

            \begin{itemize}
                \item $L(\emptyset) = \emptyset$
                \item $L(a) = \{a\} \; (a \in A)$
                \item $L(R_1 \cup R_2) = L(R_1) \cup L(R_2)$
                \item $L(R_1 R_2) = L(R_1) \cdot L(R_2)$
                \item $L(R^*) = L(R)^*$
            \end{itemize}

            \newpage

            Jede formale Sprache, die sich durch einen regulären Ausdruck (mithilfe endlich vieler Vereinigungen $\cup$,
            Konkatenationen $\cdot$ und Kleene-Sterne $^*$) beschreiben lässt,
            heißt per Definition eine reguläre Sprache.

            \emptyline

            Reguläre Ausdrücke besitzen kein unbegrenztes \afz{Gedächtnis} (sie können nicht beliebig weit zählen)
            und können daher keine Abhängigkeiten zwischen beliebig langen Teilwörtern erzwingen,
            weshalb Sprachen wie $L = \{\text{a}^n\text{b}^n : n \in \N_+\}$ nicht regulär sind. 

        \subsection{Induktive Definitionen}

            Die Definition der n-ten Potenz einer Sprache $L$ können wir mit einer induktiven Definition schreiben:

            \vspace{-1em}
            \begin{align*}
                \forall n \in \N_0: L^n =
                \begin{cases}
                    \{\varepsilon\} &\If{n = 0} \\
                    L \cdot L^{n-1} &\If{n \geq 1}
                \end{cases}
            \end{align*}

            Wir definieren $L^n$ wenn $L^{n-1}$ schon definiert ist.
            Das ist eine induktive Definition, aka. rekursive Definition.
            $L^0 = \{\varepsilon\}$ heißt die Verankerung.

            \emptyline

            Sei $A$ ein beliebiges Alphabet.
            Wir liegen zwei weitere Beispiele vor:

            \vspace{-1em}
            \begin{align*}
                \forall w \in A^*\ \forall a \in A&: |w|_a =
                \begin{cases}
                    0          &\If{w = \varepsilon} \\
                    1 + |w'|_a &\If{\exists w' \in A^*: w = aw'} \\
                    |w'|_a     &\If{w' \in A^* \wedge \exists x \in A \setminus \{a\}: w = xw'}
                \end{cases} \\
                \forall w \in A^*&: \overline{w} =
                \begin{cases}
                    w              &\If{|w| \leq 1} \\
                    \overline{w'}x &\If{\exists x \in A\ \exists w' \in A^+: w = xw'}
                \end{cases}
            \end{align*}

            Wir können mit Induktiver Definition festlegen, wie ein korrekter Klammerausdruck aussieht.
            Ein Wort $w \in A^*$ ist ein korrekter Klammerausdruck, wenn:

            \begin{itemize}
                \item $w = \varepsilon$ oder
                \item $w = w_1 w_2$ für korrekte Klammerausdrücke $w_1$ und $w_2$ oder
                \item $w = (w_1)$ für korrekten Klammerausdruck $w_1$
            \end{itemize}

            Als Beispiel können wir $\texttt{(())()()} \in L$ in einer Baumdiagramm hinschreiben:

            \begin{center}
                \scalebox{0.8}{
                    \begin{tikzpicture}[
                        level distance=18mm,
                        every node/.style={font=\large},
                        lnode/.style={
                            rectangle,
                            fill=blue!10,
                            inner sep=4.5pt,
                            minimum height=20pt
                        },
                        level 1/.style={sibling distance=72mm},
                        level 2/.style={sibling distance=16mm},
                        level 3/.style={sibling distance=16mm},
                        level 4/.style={sibling distance=16mm},
                        edge from parent/.style={draw=black, thin}
                    ]
                        \node[lnode] {$\textcolor{blue}{\texttt{(())()()}} \in L$}
                            child {
                                node[lnode] {$\textcolor{blue}{\texttt{(())}} \in L$}
                                child { node {$\textcolor{blue}{\texttt{(}}$} }
                                child {
                                    node[lnode] {$\textcolor{blue}{\texttt{()}} \in L$}
                                    child { node {$\textcolor{blue}{\texttt{(}}$} }
                                    child {
                                        node[lnode] {$\varepsilon \in L$}
                                        child { node {$\varepsilon$} }
                                    }
                                    child { node {$\textcolor{blue}{\texttt{)}}$} }
                                }
                                child { node {$\textcolor{blue}{\texttt{)}}$} }
                            }
                            child [sibling distance=72mm] {
                                node[lnode] {$\textcolor{blue}{\texttt{(())}} \in L$}
                                child [sibling distance=48mm] {
                                    node[lnode] {$\textcolor{blue}{\texttt{()}} \in L$}
                                    child [sibling distance=16mm] { node {$\textcolor{blue}{\texttt{(}}$} }
                                    child [sibling distance=16mm] {
                                        node[lnode] {$\varepsilon \in L$}
                                        child { node {$\varepsilon$} }
                                    }
                                    child [sibling distance=16mm] { node {$\textcolor{blue}{\texttt{)}}$} }
                                }
                                child [sibling distance=48mm] {
                                    node[lnode] {$\textcolor{blue}{\texttt{()}} \in L$}
                                    child [sibling distance=16mm] { node {$\textcolor{blue}{\texttt{(}}$} }
                                    child [sibling distance=16mm] {
                                        node[lnode] {$\varepsilon \in L$}
                                        child { node {$\varepsilon$} }
                                    }
                                    child [sibling distance=16mm] { node {$\textcolor{blue}{\texttt{)}}$} }
                                }
                            };
                    \end{tikzpicture}
                }
            \end{center}

            \newpage

            Es gibt eine äquivalente Charakterisierung:
            Für jedes Wort $w \in \{(, )\}^*$ gilt:
            $w$ ist ein korrekter Klammerausdruck is äquivalent zu $|w|_( = |w|_)$ und $|v|_( \geq |v|_)$ für jedes Präfix $v$ von $w$.

            \emptyline

            Der Beweis ist folgende:

            \emptyline

            \begin{adjustwidth}{2em}{0pt}

                Wir beweisen erst die Hinrichtung über die Länge $|w|$ von $w$.

                \emptyline

                IA ($|w| = 0$): Dann gilt $w = \varepsilon$, also auch $v = w = \varepsilon$ und $|v|_( = |v|_) = 0$.

                \emptyline

                IS ($|w| \geq 1$): Dann gilt $w = w_1w_2$ oder $w = (w_1)$ mit $w_1w_2$ korrekte Klammerausdrücke.

                \begin{enumerate}
                    \item $w = w_1w_2$ mit $|w_1|, |w_2| < |w|$. \\
                          Nach IV gilt $|w_1|_( = |w_1|_)$ und $|w_2|_( = |w_2|_)$ und daher $|w|_( = |w|_)$. \\
                          Seien $v_1$, $v_2$ Präfixe von $w_1$ und $w_2$, und $v_1$, $v_2$ bilden zusammen $v$.
                          Es kann sein, dass $v$ ganz in $w_1$ liegt, was heißt $v_1 = v$ und $v_2 = \varepsilon$.
                          Oder $v$ liegt nicht nur in $w_1$, sogar auch in $w_2$ ($v_1 = w_1$).
                          Nach IV gilt $|v_1|_( \geq |v_1|_)$ und $|v_2|_( \geq |v_2|_)$ und daher $|v|_( \geq |v|_)$.
                    \item $w = (w_1)$ mit $|w_1| = |w| - 2 < |w|$. \\
                          Nach IV gilt $|w_1|_( = |w_1|_)$ und daher $|w|_( = |w|_)$. \\
                          Sei $v_1$ Präfix von $w_1$.
                          Dann ist $v$ entweder \afz{$($}, \afz{$(v_1$} oder \afz{$(w_1)$}. \\
                          Nach IV gilt $|v_1|_( \geq |v_1|_)$ und daher $|v|_( \geq |v|_)$.
                \end{enumerate}

                Dann beweisen wir die Rückrichtung über $|w|$.

                \emptyline

                IA ($|w| = 0$): Dann gilt $w = \varepsilon$ und $w$ ist ein korrekter Klammerausdruck.

                \emptyline

                IS ($|w| \geq 1$): Dann ist das erste Zeichen von $w$ ein $($.

                \begin{enumerate}
                    \item Es gibt ein echtes Präfix $u$ von $w$ mit $|u|_( = |u|_)$. \\
                          Sei $w = uz$ mit $|z|_( = |z|_)$ und $|v|_( \geq |v|_)$ für jedes Präfix $v$ von $u$ oder $z$. \\
                          Mit $|u|$, $|z| < |w|$ gilt nach IV, dass $u$, $z$ korrekte Klammerausdrücke sind.
                          Dann ist auch $w = uz$ ein korrekter Klammerausdruck.
                    \item Jedes echte Präfix $u$ von $w$ hat $|u|_( > |u|_)$. \\
                          Zerlege $w$ als $w = (z)$.
                          Dann gilt $|z|_( = |z|_)$ und $|v|_( \geq |v|_)$ für jedes Präfix $v$ von $z$. \\
                          Mit $|z| < |w|$ gilt nach IV, dass $z$ ein korrekter Klammerausdruck ist.
                          Dann ist auch $w = (z)$ ein korrekter Klammerausdruck.
                \end{enumerate}

            \end{adjustwidth}

\end{document}
